\begin{proof}[Proof of \cref{obj:thm:observed-lower}]
Fix public constants \(K=(p_{\min},c_{\mathrm{lo}},c_{\mathrm{hi}})\), \(\beta\in(0,1]\), \(\kappa\in[0,2]\), and \(\tau>0\), under the envelope calibration in the statement. All constants chosen below may depend on \(K,\beta,\kappa,\tau\) and the fixed packet construction.

\begin{enumerate}
\item \textbf{Packet constants.}
For \(\kappa=0\), use the symmetric interior construction in \cref{obj:def:packet-handle}; for \(\kappa>0\), use the endpoint construction. In the latter case, its Jacobi parameter satisfies
\[
b=\frac{\kappa-1}{2}>-\frac12,
\]
and the remaining Jacobi parameters equal \(r=4\). The localization estimate is supplied by \citep[Theorem 2.4, equation (2.14)]{PetrushevXu2005}, and the local increment estimate by \citep[Theorem 2.2, equation (2.4)]{https://people.math.sc.edu/pencho/Publications/kpx-06-28-06-web.pdf}. The orthogonality and positive squared-norm formula are supplied by \citep[Table 18.3.1, Jacobi row and caption]{https://dlmf.nist.gov/18.3}, and the required Gamma-ratio asymptotic by \citep[equation (5.11.12), restricted to the positive real axis]{https://dlmf.nist.gov/5.11.E12}.

Thus \cref{obj:thm:weighted-packet-construction} supplies constants
\[
C_{\mathrm{packet}}>0,\qquad c_p>0
\]
uniformly for \(m\ge4\), including positive normalizing values, evenness, continuous differentiability of the zero extension, support in \([-1,1]\), and
\[
\psi_m(0)=1,\qquad
\|\psi_m\|_\infty\le C_{\mathrm{packet}},\qquad
\|\psi_m'\|_\infty\le C_{\mathrm{packet}}m,
\]
\[
\int_{-1}^{1}|u|^\kappa|\psi_m(u)|\,du
\le C_{\mathrm{packet}}m^{-\kappa-1},
\qquad
\int_{-1}^{1}|u|^\kappa\psi_m(u)^2\,du
\le C_{\mathrm{packet}}m^{-\kappa-1},
\]
\[
\int_{-1}^{1}|u|^\kappa u^j\psi_m(u)\,du=0
\qquad(0\le j<c_pm).
\]
% lean: CausalSmith.Stat.NoisydoseWeakdesignTransition.lower_regime_witnesses

\item \textbf{Common calibration.}
Define the cancellation constant, degree threshold, intermediate multiplier, and compact radius by
\[
\gamma:=\frac{c_p}{2},\qquad
M_{\mathrm{deg}}:=\left\lceil\max\left\{4,\frac{2}{c_p}\right\}\right\rceil,
\qquad
C_{\mathrm I}:=M_{\mathrm{deg}},\qquad
\ell_{\mathrm M}:=\frac18.
\]
Then
\[
M_{\mathrm{deg}}\ge4,\qquad
c_pM_{\mathrm{deg}}\ge2,\qquad
C_{\mathrm I}\ge4,\qquad c_pC_{\mathrm I}\ge2.
\]
Choose the intermediate radius constant explicitly as
\[
a_{\mathrm I}:=
\min\left\{
1,\frac{1}{4\sqrt{C_{\mathrm I}+1}},
\sqrt{\min\left\{\frac{\gamma}{2},
\gamma\exp(-1/\gamma-1)\right\}}
\right\}.
\]
Every entry in this minimum is positive, and hence
\[
0<a_{\mathrm I}\le1,\qquad
a_{\mathrm I}\sqrt{C_{\mathrm I}+1}\le\frac14,\qquad
a_{\mathrm I}^2\le\frac{\gamma}{2},\qquad
\frac{\exp(1)a_{\mathrm I}^2}{\gamma}\le\exp(-1/\gamma).
\]
For the compact regime, set
\[
C_{\mathrm M}:=
\max\left\{
1,\frac8\gamma,\frac{2\exp(2)\ell_{\mathrm M}^2}{\gamma}
\right\}.
\]
Consequently,
\[
C_{\mathrm M}\ge1,\qquad
\frac{\gamma C_{\mathrm M}}{\exp(1)\ell_{\mathrm M}^2}
\ge2\exp(1),\qquad
\gamma C_{\mathrm M}\ge8.
\]

Define the design coefficient and three comparison constants by
\[
c_\kappa:=(\kappa+1)2^\kappa,\qquad
B_{\mathrm D}:=8c_\kappa,\qquad
B_{\mathrm I}:=4\exp(1/2)(c_\kappa C_{\mathrm{packet}})^2,\qquad
B_{\mathrm G}:=4\exp(1/2)(2c_\kappa)^2,
\]
and set
\[
B_{\mathrm{amp}}:=\max\{B_{\mathrm D},2B_{\mathrm I}\},\qquad
\varepsilon:=
\min\left\{
\frac18,\frac{1}{8C_{\mathrm{packet}}},
\sqrt{\frac{\log(1+\tau^2)}{B_{\mathrm{amp}}}}
\right\}.
\]
These constants are positive. The amplitude satisfies
\[
\varepsilon\le\frac18,\qquad
4\varepsilon\le1,\qquad
\varepsilon C_{\mathrm{packet}}\le\frac18,
\]
\[
B_{\mathrm D}\varepsilon^2\le\log(1+\tau^2),
\qquad
2B_{\mathrm I}\varepsilon^2\le\log(1+\tau^2).
\]
Finally, choose
\[
C:=\varepsilon^2\max\{B_{\mathrm G},B_{\mathrm I}\},
\qquad
c:=2\varepsilon\min\left\{
1,
\left(\frac{a_{\mathrm I}}{\sqrt{C_{\mathrm I}+1}}\right)^\beta,
\left(\frac{\ell_{\mathrm M}}{2C_{\mathrm M}}\right)^\beta
\right\}.
\]
Both are positive, and these definitions give the comparison constants required in each regime.
% lean: CausalSmith.Stat.NoisydoseWeakdesignTransition.lower_common_amplitude_constant

\item \textbf{Uniform sample cutoff.}
Use the effective exponent and sample logarithm from \cref{obj:def:frontier-rate}, and define
\[
d:=2\beta+\kappa+1>0,\qquad
L_n:=\log(\exp(1)n),\qquad h_0:=n^{-1/d},
\]
\[
n_{\mathrm D}:=\max\{2,\lceil4^d\rceil\},\qquad
n_{\log}:=\max\{2,\lceil\exp(M_{\mathrm{deg}}^2)\rceil\},
\qquad
n_0:=\max\{n_{\mathrm D},n_{\log}\}.
\]
For every \(n\ge n_0\),
\[
0<h_0\le(4^d)^{-1/d}=\frac14,\qquad
nh_0^d=1,\qquad
M_{\mathrm{deg}}^2\le L_n.
\]
The last inequality follows by taking logarithms of
\[
\exp(M_{\mathrm{deg}}^2)\le n\le\exp(1)n.
\]
For every \(\sigma\in[0,1/4]\), also
\[
1\le\log(\exp(1)+n\sigma^d)\le L_n.
\]
Indeed, \(0\le\sigma^d\le1\), while \(n\ge2\) and \(\exp(1)\ge2\) give
\[
\exp(1)\le\exp(1)+n\sigma^d
\le\exp(1)+n\le\exp(1)n.
\]
Fix henceforth an arbitrary public space from \cref{obj:def:potential-path-space}, \(n\ge n_0\), and \(\sigma\in[0,1/4]\).
% lean: CausalSmith.Stat.NoisydoseWeakdesignTransition.lower_log_sample_cutoff

\item \textbf{Witness laws and observed densities.}
The common centered-dose density is
\[
g(t):=c_\kappa|t|^\kappa\mathbf1\{|t|\le1/2\},
\qquad t\in\mathbb R.
\]
We use the convention \(|t|^0=1\), including at \(t=0\). Its normalization follows from
\[
\int_{-1/2}^{1/2}|t|^\kappa\,dt
=\frac{2^{-\kappa}}{\kappa+1},
\qquad
\int_{\mathbb R}g(t)\,dt=1.
\]
Take mutually independent seed coordinates with
\[
P(X=0)=P(X=1)=\frac12,\qquad
\mathcal L(A-a_0)(dt)=g(t)\,dt,\qquad
Z\sim N(0,1),\qquad U\sim\operatorname{Unif}[0,1].
\]
Here the prescribed center is \(a_0=1/2\). For any continuous profile
\(\mu:[0,1]\to[1/8,7/8]\), define its structural law by the measurable pushforward
\[
P_\mu:=
\mathcal L\bigl(X,A,Z,\iota(\mu,U),\iota(\mu,U)(A),U\bigr).
\]
Thus, under \(P_\mu\),
\[
Y(\cdot)=\iota(\mu,U),\qquad
Y(a)=\mathbf1\{U\le\mu(a)\},\qquad
Y=Y(A),\qquad W=A+\sigma Z.
\]
The measurability requirements follow from the embedding and joint evaluation in \cref{obj:def:potential-path-space}. Uniformity of the independent rank gives
\[
E_{P_\mu}[Y(a)\mid X=x]=\mu(a),
\qquad
\theta(P_\mu)=\sum_{x\in\{0,1\}}\frac12\mu(a_0)=\mu(a_0).
\]

For a continuous compactly supported perturbation, write
\[
\mu_\pm(a):=\frac12\pm\delta(a-a_0),\qquad
\mu_\star(a):=\frac12,\qquad
P_\pm:=P_{\mu_\pm},\qquad P_\star:=P_{\mu_\star},
\]
whenever the profiles have the displayed range. Define
\[
Q_\pm:=\mathcal L_{P_\pm}(X,W,Y),\qquad
Q_\star:=\mathcal L_{P_\star}(X,W,Y),\qquad
v_\delta(t):=g(t)\delta(t).
\]
For \(\sigma>0\), define
\[
\phi_\sigma(w):=
\frac{1}{\sqrt{2\pi}\sigma}
\exp\left(-\frac{w^2}{2\sigma^2}\right),
\qquad
p_V:=g*\phi_\sigma,\qquad f_\delta:=v_\delta*\phi_\sigma.
\]
For \(\sigma=0\), define instead
\[
p_V:=g,\qquad f_\delta:=v_\delta.
\]
Conditioning on the latent dose and integrating the uniform rank gives the densities
\[
q_\pm(x,w,y)
=
p_x\left\{\frac{p_V(w-a_0)}2
\pm(2y-1)f_\delta(w-a_0)\right\},
\qquad
q_\star(x,w,y)=\frac{p_xp_V(w-a_0)}2,
\]
relative to counting measure in \(x\in\{0,1\}\), \(y\in\{0,1\}\), and Lebesgue measure in \(w\in\mathbb R\), with \(p_x=1/2\).

The marked probabilities belong to \([0,1]\), so
\[
0\le Q_\pm\le2Q_\star,\qquad Q_\pm\ll Q_\star.
\]
For the likelihood ratios
\[
\Lambda_\pm:=\frac{dQ_\pm}{dQ_\star},
\]
this implies
\[
0\le\Lambda_\pm\le2\quad Q_\star\text{-almost surely},
\qquad
\int(\Lambda_\pm-1)^2\,dQ_\star\le1.
\]
We use
\[
\chi^2(Q_\pm,Q_\star):=
\int(\Lambda_\pm-1)^2\,dQ_\star,
\qquad
\chi^2(P_\pm,P_\star):=\chi^2(Q_\pm,Q_\star),
\]
in agreement with the theorem's observed-law convention. Summing the squared likelihood deviations over the two marks and strata, and translating \(w\) by \(a_0\), yields
\[
\chi^2(Q_\pm,Q_\star)
=
4\int_{\{p_V>0\}}\frac{f_\delta(w)^2}{p_V(w)}\,dw.
\]
At \(\sigma=0\), both marked densities vanish wherever \(g=0\), and this identity reduces to
\[
\chi^2(Q_\pm,Q_\star)
=4\int_{\mathbb R}g(t)\delta(t)^2\,dt.
\]
% lean: CausalSmith.Stat.NoisydoseWeakdesignTransition.observed_marked_chiSq_identity

\item \textbf{Gaussian moment comparison.}
Suppose \(\sigma>0\), and let the support radius and cancellation cutoff satisfy
\[
\ell>0,\qquad J\in\mathbb N_0,\qquad
\operatorname{supp}(v_\delta)\subseteq[-\ell,\ell],\qquad
\int_{\mathbb R}t^jv_\delta(t)\,dt=0\quad(0\le j<J).
\]
Define the absolute mass and moment tail by
\[
\|v_\delta\|_1:=\int_{\mathbb R}|v_\delta(t)|\,dt,\qquad
\mathcal T_{\sigma,\ell,J}:=
\sum_{j=J}^{\infty}\frac{(\ell^2/\sigma^2)^j}{j!}.
\]

First consider \(w\ge0\). Restricting the convolution defining \(p_V(w)\) to \(0\le t\le\sigma\), which lies in the latent support, gives
\[
\begin{aligned}
p_V(w)
&=\phi_\sigma(w)\int_{\mathbb R}
g(t)\exp\left(\frac{wt}{\sigma^2}-\frac{t^2}{2\sigma^2}\right)\,dt\\
&\ge
c_\kappa\exp(-1/2)\phi_\sigma(w)\int_0^\sigma t^\kappa\,dt\\
&=
2^\kappa\exp(-1/2)\sigma^{\kappa+1}\phi_\sigma(w)\\
&\ge\exp(-1/2)\sigma^{\kappa+1}\phi_\sigma(w)>0.
\end{aligned}
\]
For \(w<0\), evenness of \(g\) and \(\phi_\sigma\) gives the same bound by replacing \(w\) with \(-w\). Hence the denominator bound holds on the entire real line.

Completing the square gives, for \(t,s,w\in\mathbb R\),
\[
\frac{\phi_\sigma(w-t)\phi_\sigma(w-s)}{\phi_\sigma(w)}
=
\exp(ts/\sigma^2)\phi_\sigma(w-t-s).
\]
Its integral in \(w\) is \(\exp(ts/\sigma^2)\). Compact support bounds the absolute triple integral by
\[
\exp(\ell^2/\sigma^2)\|v_\delta\|_1^2<\infty.
\]
Absolute Fubini therefore yields
\[
\int_{\mathbb R}\frac{f_\delta(w)^2}{\phi_\sigma(w)}\,dw
=
\iint_{\mathbb R^2}
v_\delta(t)v_\delta(s)\exp(ts/\sigma^2)\,dt\,ds.
\]
For each \(j\in\mathbb N_0\), the absolute integral of the \(j\)-th exponential-series term is bounded by
\[
\iint_{\mathbb R^2}
\frac{|t^jv_\delta(t)s^jv_\delta(s)|}{\sigma^{2j}j!}\,dt\,ds
\le
\|v_\delta\|_1^2\frac{(\ell^2/\sigma^2)^j}{j!}.
\]
The sum of these bounds is finite. Termwise integration and the cancellations give
\[
\begin{aligned}
\int_{\mathbb R}\frac{f_\delta(w)^2}{\phi_\sigma(w)}\,dw
&=
\sum_{j=0}^{\infty}
\frac{\left(\int_{\mathbb R}t^jv_\delta(t)\,dt\right)^2}
{\sigma^{2j}j!}\\
&=
\sum_{j=J}^{\infty}
\frac{\left(\int_{\mathbb R}t^jv_\delta(t)\,dt\right)^2}
{\sigma^{2j}j!}\\
&\le \|v_\delta\|_1^2\mathcal T_{\sigma,\ell,J},
\end{aligned}
\]
where the last inequality uses
\[
\left|\int_{\mathbb R}t^jv_\delta(t)\,dt\right|
\le\ell^j\|v_\delta\|_1.
\]
Combining this with the observed-density identity and the denominator bound proves
\[
\chi^2(Q_\pm,Q_\star)
\le
4\exp(1/2)\sigma^{-\kappa-1}
\|v_\delta\|_1^2\mathcal T_{\sigma,\ell,J}.
\]
% lean: CausalSmith.Stat.NoisydoseWeakdesignTransition.observed_perturbation_chiSq_moment_bound

\item \textbf{Direct profiles.}
Suppose first that \(\sigma\le h_0\). Set
\[
h=\ell=h_0,\qquad m=0,\qquad J=0,
\]
and define, for all \(t\in\mathbb R\),
\[
\delta(t):=
\varepsilon h_0^\beta
\left(1-\left(\frac{t}{h_0}\right)^2\right)^2
\mathbf1\{|t|\le h_0\}.
\]
For the increment calculation, define the clipped coordinate
\[
\zeta_{\mathrm D}(t):=\min\{|t/h_0|,1\}.
\]
Then
\[
\delta(t)=\varepsilon h_0^\beta
(1-\zeta_{\mathrm D}(t)^2)^2,
\qquad
0\le\zeta_{\mathrm D}(t)\le1,
\qquad
|\zeta_{\mathrm D}(s)-\zeta_{\mathrm D}(t)|
\le\frac{|s-t|}{h_0}.
\]
Factoring differences of squares twice gives
\[
\begin{aligned}
\left|(1-\zeta_{\mathrm D}(s)^2)^2
-(1-\zeta_{\mathrm D}(t)^2)^2\right|
&\le2|\zeta_{\mathrm D}(s)^2-\zeta_{\mathrm D}(t)^2|\\
&\le4|\zeta_{\mathrm D}(s)-\zeta_{\mathrm D}(t)|\\
&\le4\frac{|s-t|}{h_0}.
\end{aligned}
\]
Thus
\[
0\le\delta(t)\le\varepsilon h_0^\beta,\qquad
|\delta(s)-\delta(t)|
\le4\varepsilon h_0^\beta\frac{|s-t|}{h_0}.
\]
If \(|s-t|\le h_0\), then
\[
\frac{|s-t|}{h_0}
\le\left(\frac{|s-t|}{h_0}\right)^\beta,
\]
so the increment is at most \(4\varepsilon|s-t|^\beta\). If \(|s-t|>h_0\), the range bound gives
\[
|\delta(s)-\delta(t)|
\le\varepsilon h_0^\beta
\le\varepsilon|s-t|^\beta.
\]
Consequently,
\[
|\delta(s)-\delta(t)|\le4\varepsilon|s-t|^\beta.
\]
The explicit formula is continuous across its support endpoints. Since \(h_0^\beta\le1\), the common calibration ensures
\[
\mu_\pm(a)\in[1/8,7/8],\qquad
|\mu_\pm(a)-\mu_\pm(a')|\le|a-a'|^\beta
\quad(a,a'\in[0,1]).
\]
At the target,
\[
|\theta(P_+)-\theta(P_-)|
=2\delta(0)=2\varepsilon h_0^\beta
=2\varepsilon\rho_{n,\sigma}
\ge c\rho_{n,\sigma},
\]
using the direct branch of \cref{obj:def:frontier-rate}.

On \([-h_0,h_0]\), the bounds on \(g\) and \(\delta\) give
\[
|v_\delta(t)|\le c_\kappa\varepsilon h_0^{\beta+\kappa}.
\]
Outside that interval \(v_\delta=0\). Integrating over its length \(2h_0\) yields
\[
\|v_\delta\|_1
\le2c_\kappa\varepsilon h_0^{\beta+\kappa+1}.
\]
Moreover,
\[
\int_{\mathbb R}g(t)\delta(t)^2\,dt
\le\varepsilon h_0^\beta\|v_\delta\|_1
\le2c_\kappa\varepsilon^2h_0^d.
\]
For positive noise, the Gaussian moment comparison with \(J=0\) therefore gives
\[
\chi^2(Q_\pm,Q_\star)
\le
B_{\mathrm G}\varepsilon^2
\sigma^{-\kappa-1}h_0^{2\beta+2\kappa+2}
\mathcal T_{\sigma,h_0,0}
\le
C\sigma^{-\kappa-1}h^{2\beta+2\kappa+2}
\mathcal T_{\sigma,\ell,J}.
\]
The moment conditions for \(J=0\) have an empty index set.
% lean: CausalSmith.Stat.NoisydoseWeakdesignTransition.direct_lower_profile_frontier_construction

\item \textbf{Direct observed testing bound.}
At \(\sigma=0\), the exact observed comparison already gives
\[
\chi^2(Q_\pm,Q_\star)
=4\int_{\mathbb R}g(t)\delta(t)^2\,dt.
\]
At \(\sigma>0\), Cauchy--Schwarz under the finite positive measure
\(g(t)\phi_\sigma(w-t)\,dt\) gives
\[
f_\delta(w)^2
\le
p_V(w)\int_{\mathbb R}
g(t)\delta(t)^2\phi_\sigma(w-t)\,dt.
\]
Dividing by the positive denominator, integrating, and using the unit Gaussian integral yields
\[
\chi^2(Q_\pm,Q_\star)
\le4\int_{\mathbb R}g(t)\delta(t)^2\,dt.
\]
Thus, for every noise value in the direct regime,
\[
n\chi^2(Q_\pm,Q_\star)
\le8c_\kappa\varepsilon^2nh_0^d
=B_{\mathrm D}\varepsilon^2
\le\log(1+\tau^2).
\]

Here is the product comparison that will also be used for the inverse regimes. Independence in \cref{obj:def:observed-experiment} and the square-integrability of the likelihood ratios imply
\[
\begin{aligned}
1+\chi^2(Q_{P_\pm}^n,Q_{P_\star}^n)
&=
E_{Q_\star^{\otimes n}}
\left[\prod_{i=1}^n\Lambda_\pm(O_i)^2\right]\\
&=
\left(E_{Q_\star}\Lambda_\pm^2\right)^n
=\{1+\chi^2(Q_\pm,Q_\star)\}^n\\
&\le\exp\{n\chi^2(Q_\pm,Q_\star)\}.
\end{aligned}
\]
Hence the sample budget implies
\[
\chi^2(Q_{P_\pm}^n,Q_{P_\star}^n)\le\tau^2.
\]
For an absolutely continuous probability law, the signed likelihood deviation has integral zero. Its positive and negative parts consequently have equal integrals, and Cauchy--Schwarz gives
\[
\operatorname{TV}(Q_{P_\pm}^n,Q_{P_\star}^n)
=
\frac12E_{Q_{P_\star}^n}
\left|\prod_{i=1}^n\Lambda_\pm(O_i)-1\right|
\le\frac12\sqrt{\chi^2(Q_{P_\pm}^n,Q_{P_\star}^n)}
\le\frac\tau2.
\]
The triangle inequality now proves
\[
\operatorname{TV}(Q_{P_+}^n,Q_{P_-}^n)\le\tau.
\]
% lean: CausalSmith.Stat.NoisydoseWeakdesignTransition.direct_profile_pair_observed_tv_bound

\item \textbf{Inverse profiles and tail estimates.}
For the remaining cases, \(\sigma>h_0>0\). Given
\[
m\in\mathbb N,\qquad m\ge M_{\mathrm{deg}},\qquad
0<\ell\le\frac14,
\]
define
\[
h:=\frac{\ell}{m},\qquad
J:=\lfloor c_pm\rfloor,\qquad
\delta(t):=\varepsilon h^\beta\psi_m(t/\ell)
\quad(t\in\mathbb R).
\]
Then \(0<h\le1/4\). Since \(c_pm\ge2\) and
\(\lfloor c_pm\rfloor>c_pm-1\),
\[
J\ge\gamma m,\qquad J\ge1.
\]
The packet bounds imply continuity and support in \([-\ell,\ell]\), together with
\[
|\delta(t)|\le\varepsilon C_{\mathrm{packet}}h^\beta.
\]
The derivative bound and the mean-value inequality give
\[
|\delta(s)-\delta(t)|
\le\varepsilon C_{\mathrm{packet}}h^\beta\frac{|s-t|}{h}.
\]
If \(|s-t|\le h\), replacing \(|s-t|/h\) by its \(\beta\)-th power proves
\[
|\delta(s)-\delta(t)|
\le\varepsilon C_{\mathrm{packet}}|s-t|^\beta.
\]
If \(|s-t|>h\), the supremum bound gives
\[
|\delta(s)-\delta(t)|
\le2\varepsilon C_{\mathrm{packet}}h^\beta
\le2\varepsilon C_{\mathrm{packet}}|s-t|^\beta.
\]
Therefore
\[
|\delta(s)-\delta(t)|
\le2\varepsilon C_{\mathrm{packet}}|s-t|^\beta.
\]
The common amplitude calibration implies
\[
\mu_\pm(a)\in[1/8,7/8],\qquad
|\mu_\pm(a)-\mu_\pm(a')|\le|a-a'|^\beta.
\]
These arguments include the Lipschitz endpoint \(\beta=1\).

For \(j<J\), we have \(j<c_pm\). Changing variables \(t=\ell u\), using \(\ell\le1/4\), and applying the packet cancellations gives
\[
\int_{\mathbb R}t^jv_\delta(t)\,dt
=
c_\kappa\varepsilon h^\beta\ell^{\kappa+j+1}
\int_{-1}^{1}|u|^\kappa u^j\psi_m(u)\,du
=0.
\]
The weighted absolute-mass bound similarly gives
\[
\begin{aligned}
\|v_\delta\|_1
&=
c_\kappa\varepsilon h^\beta\ell^{\kappa+1}
\int_{-1}^{1}|u|^\kappa|\psi_m(u)|\,du\\
&\le
c_\kappa\varepsilon C_{\mathrm{packet}}
h^\beta\ell^{\kappa+1}m^{-\kappa-1}\\
&=
c_\kappa\varepsilon C_{\mathrm{packet}}h^{\beta+\kappa+1}.
\end{aligned}
\]
Thus the Gaussian moment comparison proves
\[
\chi^2(Q_\pm,Q_\star)
\le
B_{\mathrm I}\varepsilon^2
\sigma^{-\kappa-1}h^{2\beta+2\kappa+2}
\mathcal T_{\sigma,\ell,J}
\le
C\sigma^{-\kappa-1}h^{2\beta+2\kappa+2}
\mathcal T_{\sigma,\ell,J}.
\]
Normalization at zero gives the exact separation
\[
|\theta(P_+)-\theta(P_-)|
=2\varepsilon h^\beta.
\]

To control the sample budget, define
\[
\lambda:=\frac{\ell^2}{\sigma^2}>0.
\]
Whenever \(J\ge1\) and \(J\ge2\lambda\), consecutive factorial-series terms have ratio at most \(1/2\). Induction therefore gives
\[
\frac{\lambda^{J+k}}{(J+k)!}
\le\frac{\lambda^J}{J!}\,2^{-k}
\quad(k\in\mathbb N_0),
\qquad
\mathcal T_{\sigma,\ell,J}
\le2\frac{\lambda^J}{J!}.
\]
The nonnegative exponential series gives
\[
\frac{J^J}{J!}
\le\sum_{k=0}^\infty\frac{J^k}{k!}
=\exp(J).
\]
Multiplication by \((\lambda/J)^J\) yields
\[
\frac{\lambda^J}{J!}
\le\left(\frac{\exp(1)\lambda}{J}\right)^J,
\qquad
\mathcal T_{\sigma,\ell,J}
\le2\left(\frac{\exp(1)\lambda}{J}\right)^J.
\]
Finally, whenever \(h\le\sigma\),
\[
\sigma^{-\kappa-1}h^{2\beta+2\kappa+2}
=
h^d(h/\sigma)^{\kappa+1}
\le h^d.
\]
% lean: CausalSmith.Stat.NoisydoseWeakdesignTransition.inverse_lower_profile_chiSq_construction

\item \textbf{Intermediate regime.}
Suppose
\[
h_0<\sigma\le L_n^{-1/2}.
\]
Define the effective sample quantity and its logarithm by
\[
N_{\mathrm I}:=n\sigma^d,\qquad
S_{\mathrm I}:=\log(\exp(1)+N_{\mathrm I}),
\]
and choose
\[
m:=\lceil C_{\mathrm I}S_{\mathrm I}\rceil,\qquad
\ell:=a_{\mathrm I}\sigma\sqrt m,\qquad
h:=\frac{\ell}{m}=\frac{a_{\mathrm I}\sigma}{\sqrt m},
\qquad
J:=\lfloor c_pm\rfloor.
\]
The logarithmic window gives \(1\le S_{\mathrm I}\le L_n\). The ceiling bounds imply
\[
C_{\mathrm I}S_{\mathrm I}\le m
<C_{\mathrm I}S_{\mathrm I}+1
\le(C_{\mathrm I}+1)S_{\mathrm I}.
\]
In particular \(m\ge M_{\mathrm{deg}}\), so \(J\ge\gamma m\) and \(J\ge1\). Moreover,
\[
\begin{aligned}
0<\ell
&\le a_{\mathrm I}\sigma
\sqrt{C_{\mathrm I}+1}\sqrt{L_n}\\
&\le a_{\mathrm I}\sqrt{C_{\mathrm I}+1}
\le\frac14,
\end{aligned}
\qquad
0<h\le\ell\le\frac14,\qquad h\le\sigma.
\]
Thus the inverse profiles just constructed are legal. Their bandwidth satisfies
\[
h=\frac{a_{\mathrm I}\sigma}{\sqrt m}
\ge
\frac{a_{\mathrm I}}{\sqrt{C_{\mathrm I}+1}}
\frac{\sigma}{\sqrt{S_{\mathrm I}}},
\]
and therefore
\[
|\theta(P_+)-\theta(P_-)|
\ge
2\varepsilon
\left(\frac{a_{\mathrm I}}{\sqrt{C_{\mathrm I}+1}}\right)^\beta
\left(\frac{\sigma}{\sqrt{S_{\mathrm I}}}\right)^\beta
\ge c\rho_{n,\sigma},
\]
using the intermediate branch of \cref{obj:def:frontier-rate}.

For the Gaussian tail parameter,
\[
\lambda=\frac{\ell^2}{\sigma^2}=a_{\mathrm I}^2m,
\qquad
2\lambda\le\gamma m\le J,
\]
and the radius calibration gives
\[
\frac{\exp(1)\lambda}{J}
\le\frac{\exp(1)a_{\mathrm I}^2}{\gamma}
\le\exp(-1/\gamma).
\]
Consequently,
\[
\mathcal T_{\sigma,\ell,J}
\le2\exp(-J/\gamma)\le2\exp(-m).
\]
Since \(m\ge S_{\mathrm I}\),
\[
N_{\mathrm I}\exp(-m)
\le
N_{\mathrm I}\exp(-S_{\mathrm I})
=
\frac{N_{\mathrm I}}{\exp(1)+N_{\mathrm I}}
\le1.
\]
Using \(h\le\sigma\) and \(d>0\), we obtain the complete sample-factor bound
\[
\begin{aligned}
n\sigma^{-\kappa-1}h^{2\beta+2\kappa+2}
\mathcal T_{\sigma,\ell,J}
&\le2n\sigma^d\exp(-m)\\
&=2N_{\mathrm I}\exp(-m)\le2.
\end{aligned}
\]
Hence
\[
n\chi^2(Q_\pm,Q_\star)
\le2B_{\mathrm I}\varepsilon^2
\le\log(1+\tau^2).
\]
The product comparison established above proves
\[
\operatorname{TV}(Q_{P_+}^n,Q_{P_-}^n)\le\tau.
\]
% lean: CausalSmith.Stat.NoisydoseWeakdesignTransition.intermediate_lower_profile_sample_budget

\item \textbf{Compact regime.}
The remaining case satisfies
\[
\sigma>h_0,\qquad \sigma>L_n^{-1/2}.
\]
Define
\[
\tau_{\mathrm M}:=\sigma^2L_n>1,\qquad
B_{\mathrm M}:=\log(\exp(1)+\tau_{\mathrm M}),
\]
and choose
\[
m:=\left\lceil\frac{C_{\mathrm M}L_n}{B_{\mathrm M}}\right\rceil,
\qquad
\ell:=\ell_{\mathrm M},\qquad
h:=\frac{\ell_{\mathrm M}}m,\qquad
J:=\lfloor c_pm\rfloor.
\]

We first establish the logarithmic bounds used for this choice. Since
\(\tau_{\mathrm M}\ge1\),
\[
\exp(1)+\tau_{\mathrm M}
\le(\exp(1)+1)\tau_{\mathrm M}.
\]
Also \(\exp(1)+1\le\exp(2)\), so
\[
B_{\mathrm M}
\le\log\tau_{\mathrm M}+\log(\exp(1)+1)
\le\log\tau_{\mathrm M}+2.
\]
Applying \(\log x\le x-1\) to \(x=\sqrt{\tau_{\mathrm M}}\) gives
\[
\log\tau_{\mathrm M}\le2\sqrt{\tau_{\mathrm M}}-2,
\qquad
B_{\mathrm M}\le2\sqrt{\tau_{\mathrm M}}.
\]
Because \(\sigma\le1/4\),
\[
0<B_{\mathrm M}
\le2\sigma\sqrt{L_n}
\le\frac{\sqrt{L_n}}2,
\qquad
\frac{L_n}{B_{\mathrm M}}\ge2\sqrt{L_n}\ge1.
\]
The ceiling bounds consequently yield
\[
\frac{C_{\mathrm M}L_n}{B_{\mathrm M}}
\le m\le\frac{2C_{\mathrm M}L_n}{B_{\mathrm M}}.
\]
Furthermore,
\[
m\ge\frac{L_n}{B_{\mathrm M}}
\ge2\sqrt{L_n}\ge M_{\mathrm{deg}},
\]
where \(M_{\mathrm{deg}}^2\le L_n\) was enforced by the sample cutoff. Thus
\[
J\ge\gamma m\ge
\frac{\gamma C_{\mathrm M}L_n}{B_{\mathrm M}},
\qquad J\ge1.
\]
The bandwidth is positive and at most \(\ell_{\mathrm M}\le1/4\). Since
\(\sigma\sqrt{L_n}\ge1\),
\[
\sigma m\ge2\sigma\sqrt{L_n}\ge2\ge\ell_{\mathrm M},
\qquad h\le\sigma.
\]
The upper degree bound gives
\[
h=\frac{\ell_{\mathrm M}}m
\ge
\frac{\ell_{\mathrm M}}{2C_{\mathrm M}}\frac{B_{\mathrm M}}{L_n}.
\]
The inverse-profile separation therefore satisfies
\[
|\theta(P_+)-\theta(P_-)|
\ge
2\varepsilon
\left(\frac{\ell_{\mathrm M}}{2C_{\mathrm M}}\right)^\beta
\left(\frac{B_{\mathrm M}}{L_n}\right)^\beta
\ge c\rho_{n,\sigma},
\]
using the compact branch of \cref{obj:def:frontier-rate}.

For the tail calculation, define
\[
\lambda:=\frac{\ell_{\mathrm M}^2}{\sigma^2}
=\frac{\ell_{\mathrm M}^2L_n}{\tau_{\mathrm M}},
\qquad
D_{\mathrm M}:=
\frac{\gamma C_{\mathrm M}}
{\exp(1)\ell_{\mathrm M}^2}\ge2\exp(1).
\]
The cancellation order and the bound \(B_{\mathrm M}\le2\sqrt{\tau_{\mathrm M}}\) imply
\[
\frac{J}{\exp(1)\lambda}
\ge
D_{\mathrm M}\frac{\tau_{\mathrm M}}{B_{\mathrm M}}
\ge\exp(1)\sqrt{\tau_{\mathrm M}}
\ge2.
\]
In particular \(J\ge2\lambda\). Taking logarithms, and using
\(\log(\exp(1)+1)\le2\), gives
\[
\log\frac{J}{\exp(1)\lambda}
\ge1+\frac12\log\tau_{\mathrm M}
\ge\frac{B_{\mathrm M}}2.
\]
Hence
\[
J\log\frac{J}{\exp(1)\lambda}
\ge
\frac{\gamma C_{\mathrm M}L_n}{B_{\mathrm M}}
\frac{B_{\mathrm M}}2
=
\frac{\gamma C_{\mathrm M}L_n}{2}
\ge4L_n.
\]
The factorial-tail estimate now gives
\[
\mathcal T_{\sigma,\ell_{\mathrm M},J}
\le
2\exp\left(-J\log\frac{J}{\exp(1)\lambda}\right)
\le2\exp(-4L_n).
\]
Using \(h\le\sigma\) for the first inequality and \(0<h\le1\) for the second, the Gaussian prefactor satisfies
\[
\sigma^{-\kappa-1}h^{2\beta+2\kappa+2}\le h^d\le1.
\]
Also \(L_n\ge0\), so
\[
n\exp(-4L_n)
\le n\exp(-L_n)
=\exp(-1)\le1.
\]
It follows that
\[
n\sigma^{-\kappa-1}h^{2\beta+2\kappa+2}
\mathcal T_{\sigma,\ell_{\mathrm M},J}\le2,
\]
and therefore
\[
n\chi^2(Q_\pm,Q_\star)
\le2B_{\mathrm I}\varepsilon^2
\le\log(1+\tau^2).
\]
The same product comparison proves
\[
\operatorname{TV}(Q_{P_+}^n,Q_{P_-}^n)\le\tau.
\]
% lean: CausalSmith.Stat.NoisydoseWeakdesignTransition.compact_lower_profile_sample_budget

\item \textbf{Structural membership and assembly.}
Each regime has supplied continuous profiles satisfying
\[
\mu_\pm(a)\in[1/8,7/8],\qquad
|\mu_\pm(a)-\mu_\pm(a')|\le|a-a'|^\beta.
\]
The reference profile satisfies the same conditions, since
\[
\mu_\star(a)=\frac12,\qquad
|\mu_\star(a)-\mu_\star(a')|=0\le|a-a'|^\beta.
\]
We now verify membership for all three pushforward laws.

The common seed gives
\[
p_0=p_1=\frac12,\qquad
A=a_0+t\in[0,1]\quad\text{almost surely}.
\]
The common conditional centered-dose density has coefficient
\(c_\kappa=(\kappa+1)2^\kappa\). The assumed calibration therefore gives
\[
c_{\mathrm{lo}}|t|^\kappa\le g(t)
\le c_{\mathrm{hi}}|t|^\kappa
\qquad(t\in[-1/2,1/2]).
\]
Because the class constants satisfy \(p_{\min}\le1/2\), the equal stratum probabilities obey \(p_{\min}\le p_x=1/2\le1-p_{\min}\) for both strata. These facts establish \cref{obj:ass:stratum-overlap,obj:ass:latent-support,obj:ass:weak-design}.

For every one of the three profiles, the conditional-mean identity derived from the uniform seed is
\[
E_{P_\mu}[Y(a)\mid X=x]=\mu(a).
\]
The established profile bounds therefore give
\cref{obj:ass:holder-mean,obj:ass:mean-range}. The chosen Gaussian marginal gives \cref{obj:ass:gaussian-channel}. Independence of the Gaussian seed from the other seed coordinates is preserved under the measurable schedule map, yielding
\[
Z\perp(X,A,Y(\cdot)),
\]
as required by \cref{obj:ass:error-independence}.

The product seed also gives
\[
\mathcal L(U\mid X=x)=\operatorname{Unif}[0,1],
\qquad U\perp A\mid X,
\]
establishing \cref{obj:ass:rank-uniform,obj:ass:rank-treatment-independence}. More explicitly, for measurable schedule and dose events,
\[
\begin{aligned}
&P_\mu\{Y(\cdot)\in\mathcal B_{\mathrm{sched}},
A\in\mathcal B_{\mathrm{dose}}\mid X=x\}\\
&\qquad=
P\{\iota(\mu,U)\in\mathcal B_{\mathrm{sched}}\}
P\{A\in\mathcal B_{\mathrm{dose}}\}.
\end{aligned}
\]
Here the events have domains
\[
\mathcal B_{\mathrm{sched}}\in\mathcal E,\qquad
\mathcal B_{\mathrm{dose}}\in\mathcal B([0,1]).
\]
This factorization proves \cref{obj:ass:conditional-unconfoundedness} for the entire random schedule.

The measurable pushforward construction carries a single full-measure event on which
\[
Y(a)=\mathbf1\{U\le\mu(a)\}
=\mathbf1\{U\le\mu_X(a)\}
\qquad\text{for every }a\in[0,1].
\]
Thus \cref{obj:ass:structural-threshold,obj:ass:binary-potential-outcomes} hold. The same identity gives \(Y(a)\in[0,1]\), and the construction gives \(Y=Y(A)\); these establish \cref{obj:ass:bounded-potential-outcomes,obj:ass:consistency}. All ten conditions in \cref{obj:def:model-class} are therefore satisfied by \(P_+,P_-,P_\star\).

For each law, take the observed product measure
\[
Q_P^n:=\bigotimes_{i=1}^n\mathcal L_P(X,W,Y).
\]
This is the experiment in \cref{obj:def:observed-experiment} and satisfies \cref{obj:ass:iid}. Finally, projection of the seed pushforward onto \((X,A)\) removes the profile-dependent coordinates, so
\[
\mathcal L_{P_+}(X,A)
=\mathcal L_{P_-}(X,A)
=\mathcal L_{P_\star}(X,A).
\]
The direct branch includes \(\sigma=0\) and \(\sigma=h_0\); among the remaining noise values, the intermediate branch includes \(\sigma=L_n^{-1/2}\), and the compact branch covers the strict upper side. Each branch has supplied the required separation, observed product-TV bound, and, for positive noise, the displayed Gaussian moment-series comparison with the same constant \(C\). The three laws use the same seed and the same regime-specific tuning data, as asserted.
% lean: CausalSmith.Stat.NoisydoseWeakdesignTransition.observed_lower
\end{enumerate}
\end{proof}
